\documentclass[11pt]{article}
\usepackage{amsmath,amssymb,amsthm,mathtools}
\usepackage{geometry}
\geometry{margin=1in}
\newtheorem{theorem}{Theorem}
\newtheorem{lemma}{Lemma}
\newtheorem{corollary}{Corollary}
\newtheorem{definition}{Definition}
\newtheorem{remark}{Remark}
\title{Step 135: Incidence-Conditioned $\Omega$-Tail Theorem}
\author{Six Birds / RH Membrane Working Notes}
\date{}
\begin{document}
\maketitle

\section{Purpose}
Steps 130--134 isolated the following obstruction.  The BPRZ shifted second moment gives a natural source-weighted Gram, but the shifted main kernel has squarefree Boolean near-null directions.  The actual Burnol/Müntz residual class is not arbitrary: it factors through a finite zeta/Müntz incidence convolution.  Deep Boolean blind modes only see high-divisibility seed coefficients.  Step 135 formulates the exact response-space estimate required to turn that observation into a lawful source record.

\section{Finite Boolean incidence model}
Let $P_y$ be a finite set of primes, $|P_y|=k$.  We identify squarefree divisors of $P_y$ with subsets $S\subseteq P_y$.  The finite zeta/incidence convolution is
\[
   (Z_y b)(T)=\sum_{S\subseteq T} b(S).
\]
Let $\mathcal B$ be a declared family of Boolean blind modes.  The upward closure of $\mathcal B$ is
\[
   \uparrow\mathcal B=\{S\subseteq P_y:\exists A\in\mathcal B,\ A\subseteq S\}.
\]
Let $P_{\uparrow\mathcal B}$ be the seed-space projection onto this upward closure and let $\Pi_{\mathcal B}$ be the response-space projection onto the blind Walsh modes.

\begin{lemma}[upward-closure exactness]
For every seed $b$,
\[
   \Pi_{\mathcal B}Z_y b=\Pi_{\mathcal B}Z_yP_{\uparrow\mathcal B}b.
\]
Thus the blind response is completely determined by the high-divisibility seed tail.
\end{lemma}

\section{The incidence-conditioned tail}
Let
\[
   B_N:Y_{B,N}\to \ell^2(2^{P_y})
\]
be the actual regularized Burnol/Müntz residual seed synthesis map, with Gram
\[
   G_{B,N}=B_N^*B_N.
\]
Let $P_{\Omega\le K}$ be an upstream-declared Heap--Soundararajan style block cutoff:
\[
   \Omega_j(S)\le K_j\qquad\text{for each block }j.
\]
Define the incidence-conditioned high-divisibility defect by
\[
   \boxed{
   \tau_{\Omega,N}^{Z}
   :=\left\|Z_y(I-P_{\Omega\le K})B_NG_{B,N}^{-1/2}\right\|.
   }
\]
This is the correct tail.  The unconditioned seed tail
\[
   \|(I-P_{\Omega\le K})B_NG_{B,N}^{-1/2}\|
\]
is not enough, because $Z_y$ may amplify high-divisibility coefficients.

\section{Incidence-conditioned suppression theorem}
\begin{theorem}[incidence-conditioned suppression]
Assume the declared blind family $\mathcal B$ is covered by the block cutoff, meaning every blind negative-prime set violates at least one block cutoff unless it lies inside the tail:
\[
   \uparrow\mathcal B\subseteq\{S:\exists j,\ \Omega_j(S)>K_j\}.
\]
Then
\[
   \Pi_{\mathcal B}Z_yB_NG_{B,N}^{-1/2}
   =\Pi_{\mathcal B}Z_y(I-P_{\Omega\le K})B_NG_{B,N}^{-1/2},
\]
and hence
\[
   \boxed{
   \left\|\Pi_{\mathcal B}Z_yB_NG_{B,N}^{-1/2}\right\|
   \le \tau_{\Omega,N}^{Z}.
   }
\]
Consequently, for the GCD-log blind residual
\[
   \Xi_{{\rm GCD},N}=R_N^*\Pi_{\mathcal B}R_N,
   \qquad R_N=Z_yB_N,
\]
one has the Loewner bound
\[
   \boxed{
   \Xi_{{\rm GCD},N}\preceq (\tau_{\Omega,N}^{Z})^2 G_{B,N}.
   }
\]
\end{theorem}

\section{Exponential moment sufficient condition}
The previous theorem is exact but still shifts the burden to $\tau_{\Omega,N}^{Z}$.  A usable sufficient condition is an exponential moment bound.

Let $\mathcal E_{j,\sigma}$ denote multiplication by $e^{\sigma\Omega_j(S)}$ on seed coefficients.  If
\[
   \left\|\mathcal E_{j,\sigma}B_NG_{B,N}^{-1/2}\right\|\le E_{j,N}(\sigma)
\]
and if the incidence tail operator has norm $C_{Z,j,N}$ on the corresponding block-tail support, then
\[
   \left\|Z_y1_{\Omega_j>K_j}B_NG_{B,N}^{-1/2}\right\|
   \le
   C_{Z,j,N}e^{-\sigma K_j}E_{j,N}(\sigma).
\]
Using a union bound over blocks gives
\[
   \boxed{
   \tau_{\Omega,N}^{Z}
   \le
   \sum_j C_{Z,j,N}e^{-\sigma_jK_j}E_{j,N}(\sigma_j).
   }
\]
The raw Boolean model has $\|Z_y\|=\varphi^k$ with $\varphi=(1+\sqrt5)/2$, so in an unweighted model one must beat the incidence amplification.  A lawful response-space theorem must either use a better weighted geometry or prove that the actual seed has tail decay strong enough to dominate this factor.

\section{Restricted BPRZ salvage}
Suppose the GCD-log main kernel has a lower frame on the nonblind response sector:
\[
   K_q\succeq \gamma_q I\quad\text{on }(I-\Pi_{\mathcal B})H,
   \qquad \gamma_q\asymp \log q.
\]
If
\[
   \left\|\Pi_{\mathcal B}R_NG_{B,N}^{-1/2}\right\|
   \le \delta_N<1,
\]
then
\[
   \boxed{
   R_N^*K_qR_N\succeq \gamma_q(1-
   \delta_N^2)G_{B,N}.
   }
\]
Thus exact suppression $\delta_N\to0$ is sufficient, but a uniform floor $\delta_N<1$ is enough for source-coercivity growth when $\gamma_q\to\infty$.

\section{Status}
Step 135 does not prove the actual Burnol/Müntz seed estimate in full generality.  It proves the exact response-space gate:
\[
   \boxed{
   \left\|Z_y(I-P_{\Omega\le K})B_NG_{B,N}^{-1/2}\right\|\le \tau_{\Omega,N}.
   }
\]
The route is accepted only if $\tau_{\Omega,N}\to0$ or $\tau_{\Omega,N}\le\tau<1$ together with a diverging source lower frame.  Otherwise $\Xi_{\rm GCD}$ remains a genuine residual requiring another source record or a scoped nonclaim.
\end{document}
